\documentclass{article}
% Official NeurIPS 2026 style, from
% https://media.neurips.cc/Conferences/NeurIPS2026/Formatting_Instructions_For_NeurIPS_2026.zip
% dblblindworkshop is the option for a double-blind workshop track, which is what
% "Who Verifies the Agents?" runs. Swap to [preprint] to see the author block while drafting.
\usepackage[dblblindworkshop]{neurips_2026}
\usepackage[utf8]{inputenc}
\usepackage[T1]{fontenc}
\usepackage{booktabs}
\usepackage{graphicx}
\usepackage{amsmath}
\usepackage{url}
\usepackage[hidelinks]{hyperref}
\usepackage{microtype}

\title{Compute Over the Graph, Do Not Show It\\
\large What an execution graph is worth to a self-repairing agent, measured against
injected faults}

\author{Fabio Rovai\thanks{Kampakis and Co Ltd, trading as The Tesseract Academy.} \and Raluca Crisan\thanks{Etiq AI. Authorship and affiliation order to be confirmed.}}
\date{Draft, \today}

\begin{document}
\maketitle

\begin{abstract}
An execution graph records what a generated pipeline did rather than what a model said it
did, and is increasingly proposed as the substrate an agent uses to repair itself. We ask
what that substrate is worth by injecting faults whose location is known by construction,
so trust decisions can be scored instead of trusted. The graph pays off through
computation over its structure, not through presentation to a model. Computing over it
works: a selector walking the dataflow to the causal root identifies the faulty boundary
in 21 of 21 contested cases where a source-order baseline manages none, and its advantage
is exactly the extent to which declaration order diverges from execution order. Showing it
to a model is more delicate than expected: a context region selected in advance is
indistinguishable from a size-matched random one, yet a protocol that withholds nested
detail and lets the reviewer request it one step at a time matches sending everything on
less context and does separate from random, so the benefit lies in the interaction rather
than the slice. The sharpest evidence is where checks are placed. Over 90 data faults,
hand-written contracts detect 26 and localise \emph{none}, firing every time on an
aggregating stage downstream of the cause, which is the downstream-blame bias the model
judges showed reproduced in a signal containing no model. Checks learned from caught
faults and placed where they entered detect 60 and localise all 60, with zero false
suspicion on held-out healthy runs once vocabulary checks are restricted to domains the
healthy population shows to be closed. Five measurement artifacts caught during this work,
each of which would have published a confident wrong number, argue that such a benchmark
is worth as much for what it invalidates as for what it reports.
\end{abstract}

\section{Introduction}

An agent that works for a long time on one objective will eventually get something wrong,
and what happens next is a trust decision. It has to say which earlier steps still stand,
which are contaminated, and where to resume. Wrong in one direction it redoes sound work.
Wrong in the other it builds on a spoiled result, which is worse and much harder to see.

The architectural answer under study is the execution graph. A scanner captures the
artifacts a generated pipeline produced and the functions that produced them, and that
record becomes the durable state the agent reasons over. The appeal is that a graph built
from execution cannot be talked out of what happened. Recent work treats provenance of this
kind as the substrate for agent trust \cite{wang2026traces} and as the basis for
recoverable execution from recorded checkpoints \cite{zhuang2026agentrewind}.

We take the architecture seriously enough to ask what it is worth, and the answer we find
is narrower and more useful than the one usually assumed. \textbf{The graph pays off
through computation over its structure, and not through presentation to a model.} Those are
different uses of the same artifact and they do not succeed together. A selector that walks
the recorded dataflow to the causal root of a set of flagged boundaries is exact where a
baseline without that relation is useless. Acceptance criteria attached to boundaries catch
faults a model reading the same evidence does not. But handing a model a graph-selected
region of context, which is the use the substrate is most often justified by, cannot be
distinguished from handing it a size-matched random region, and accumulating history costs
two orders of magnitude more context for no gain.

Answering the question at all required a measurement the system does not have. Trust labels
here are produced by one model reviewing another, so nothing in the loop holds an
independent answer; when a reviewer calls a boundary trusted, no artifact can contradict it.
There is good reason not to take that on faith, since judges prefer text they find familiar
\cite{wataoka2024selfpreference} and a judge's agreement with itself can be high while its
agreement with the truth is not \cite{norman2026reliability}, and a verifier checking only
what it can express admits confident false positives \cite{helff2026gaming}. We supply the missing answer
by injecting a fault into a named boundary through the same scope machinery a repair must
pass, so the poisoned boundary is confined by mechanism rather than by assertion, and what
it contaminates is read from the pipeline's own dataflow. We are not aware of prior work
scoring an agent's own boundary trust labels against faults injected into its own pipeline
with known ground truth.

The paper's second contribution is methodological and, we think, the more transferable one.
Five times during this work a measurement produced a confident number that was an artifact
of how we had built the experiment rather than a property of the system: a static-code arm
that scored 75\,\% because our faults were source mutations and collapsed to 0\,\% once
they moved into recorded data; two models that returned a fixed answer regardless of input
while appearing to score respectably; and a selection result we reported as significant on a
third of its eventual sample and had to retract. Each would have survived into a published
table. We report them alongside the results because a benchmark for verifying agents that
cannot catch its own defects has no standing to measure anyone else's.
Extended related work appears in Appendix~\ref{app:related}.

\section{The system under test}

The pipeline is a marketing research agent. A model authors Python that discovers market
demand from public sources, a scanner executes it and captures the values crossing each
function boundary, a reviewing session judges each boundary, and a repairing session
returns a change confined to one boundary.

Three properties matter. The generated pipeline declares its own review boundaries as named
functions, and the review and repair machinery anchors to those declarations. Stage inputs
and outputs pass by argument and return value rather than shared state, which makes a stage
replayable against its recorded inputs. And a run records a network cassette, so a pipeline
that fetches third-party sources replays against exactly the responses the original run
saw. The cassette is what makes this work deterministic, and later what lets us move faults
out of the source entirely.

\section{What the published comparison can and cannot settle}
\label{sec:cannot}



The published comparison ran four arms over the same frozen pipeline, input and recorded
corpus, differing in what the reviewer sees: stage metadata, accumulated history, the
complete graph, or a selected part of it (table in Appendix~\ref{app:published}). It is
carefully controlled: the same frozen pipeline, input
and recorded corpus across arms, fresh sessions, no cross-arm visibility. Three properties
nonetheless prevent it settling which regime is better, and they determine what this paper
had to build.

Sample size is the first. One run per arm and three repairs within each makes every
quantity a count out of three or a single ordinal, and the recommended arm ties the
accumulated-history baseline on effective repairs. Second, the arms were not adjudicated
identically: the archived record states the first common full-graph judge exceeded the
model's 1{,}048{,}576 character input limit on the history branch, and those runs were
rejudged with isolated graph-selected judges, a failure that also removed that arm's
timing. Third, and most usefully, the function selecting which boundary to repair differs
across arms. The graph arms filter flagged units to causal roots using upstream
identifiers; the baselines rotate with no causal step. The comparison is presented as one
of context regimes but the arms also differ in how they choose a target. Rather than
control that away we make it the treatment, and it produces the paper's largest effect.

\section{Injecting a fault whose location is known}
\label{sec:method}

The missing ingredient is an answer. If we know which boundary is wrong, a trust
signal stops being something the system relies on and becomes something we can score.

We obtain that answer by injecting the fault ourselves. A fault is a small
transformation of one boundary function's source. It is applied by locating the
function, rewriting its abstract syntax tree, and splicing the result back through
the same scope validator a repair must pass. That last detail is what makes the
ground truth trustworthy rather than asserted: an injected fault is confined to one
boundary by the same mechanism that confines a repair, so if the confinement is
broken the injection fails rather than silently spreading.

Contamination is read from the pipeline rather than declared by us. A boundary is
downstream of the injected one when the injected function's result reaches its
arguments, either directly as a nested call or through a module-level name bound to
that result. Ground truth is therefore a pair: the injected boundary, and the
ordered list of boundaries its output reaches.

Two guards keep the corpus honest. An injection that leaves the source unchanged is
refused, because a fault nothing can detect would enter the corpus and depress every
signal's score equally, making every arm look worse without distinguishing them. And
a test asserts that every fault in the catalogue still changes something on a
representative pipeline; when we added it, it caught one that did not.

\subsection{The fault catalogue}

The catalogue separates faults observed in this system's own operational logs from
synthetic ones, because the two carry different weight. Three are observed. A field
disappears from a stage's output after upstream source reads failed, which this
system's logs record alongside HTTP 429, 403 and 500 responses from four source
registries. A collected sequence is shortened while the stage still reports success.
An identifier is replaced by a plausible placeholder, which is what a schema-bound
worker actually returned before a reviewing model discarded its output. Two are
synthetic: a comparison operator is inverted, and a ranking is reversed while every
value in it stays present and valid.

An invariant here is an acceptance criterion attached to a boundary and checked
deterministically against the value it produced; a concrete example is given in
Appendix~\ref{app:invariant}.

\subsection{What a trust signal is scored on}

A signal labels every boundary trusted or suspect. Scoring compares those labels
with the injected fault and reports four things.

\emph{False trust} is the injected boundary being labelled trusted. This is the
safety failure, and a signal that labels everything trusted scores perfectly on cost
and fails completely here.

\emph{False suspicion} is a boundary that is neither the injected one nor
contaminated by it being flagged. This is the cost, paid in redone work.

\emph{Localisation error} is signed, and the sign carries the finding. The resume
point is the earliest suspect boundary, which is where a run acting on these labels
would restart. A negative error means it restarts upstream of the fault and redoes
sound work. A positive error means it restarts past the fault and keeps output the
fault has already contaminated. Averaging the two would hide the second failure
inside the first, so we report them separately.

\emph{Shipped contaminated output} records whether the positive case occurred.

\subsection{Arms}

Four signals are in scope. Structural schema validation checks a stage's output
against its declared JSON schema. Content invariants check criteria the boundary
declared in advance: required fields, minimum counts, value patterns and ordering.
Both are deterministic and cost nothing per fault, which is what allows the corpus to
be large. The third arm is the reviewing model itself, which is the only arm that
costs tokens and the only one whose labels vary between runs on the same input.

\section{Results}
\label{sec:results}

The results divide along the paper's claim. Sections~\ref{sec:compute1} to
\ref{sec:compute2} concern computing over the graph's structure, where it pays. The final
part concerns showing the graph to a model, where it does not.

\subsection{Detection: what declared signals catch}

A first corpus crashed on seventeen of twenty-eight faults, so execution is itself a strong
free detector; but a fault that halts the run never reaches a trust decision, and real
generated pipelines read inputs through accessors with defaults. Everything below concerns
faults such a pipeline survives silently: 104 cases across two pipelines, the cross product
of five fault classes, the boundaries and the applicable parameters, most from classes
observed in this system's own logs.

Structural schema validation leaves 78\,\% undetected, declared invariants 27\,\%, and a
model judge 18\,\%, none of them raising a single false suspicion. The intervals for the
latter two overlap and an exact McNemar test does not separate them ($p=0.16$), so we make
no claim that either is better; what the sample supports is that both are significantly
worse than their union, which leaves 7\,\% ($p<0.0001$ and $p=0.0005$). The failures are
disjoint. Invariants catch twelve faults the judge misses, every one a matter of counting
or shape; the judge catches twenty-one the invariants miss, eighteen of them substituted
identifiers including every plausible one that defeated the pattern check. Schema
validation contributes nothing either does not, detecting only the fault class that
changes a record's shape.

\subsection{What the graph is actually for}
\label{sec:compute1}

Detection tells a run something is wrong. It does not say what to repair. When several
boundaries are flagged, something must choose, and that choice is where the graph does work
no cheaper signal can. We measure a selector that walks recorded dataflow to the causal
root against one that walks the declared boundary list in order.

The comparison is only meaningful when those two orders differ. In a pipeline whose
functions are declared in the order they run, the causal root of any flagged set is simply
the first of them, and the selectors cannot disagree. Models writing code do not reliably
declare in execution order, so we test both conditions across four shapes: a deep chain, a
wide fan-out, a diamond that splits and reconverges, and a short chain.

\begin{table}[t]
\centering
\begin{tabular}{llrrr}
\toprule
Shape & Declaration order & Contested & Causal root & Source order \\
\midrule
chain-8    & dataflow  & 16 & 100\,\% & 100\,\% \\
chain-8    & scrambled & 16 & 100\,\% & 38\,\% \\
fan-6      & scrambled &  2 & 100\,\% & 0\,\% \\
diamond-4  & scrambled & 14 & 100\,\% & 14\,\% \\
chain-3    & scrambled &  6 & 100\,\% & 0\,\% \\
\midrule
all        & dataflow  & 38 & 100\,\% & 100\,\% \\
all        & scrambled & 38 & \textbf{100\,\%} & \textbf{21\,\%} \\
\bottomrule
\end{tabular}
\caption{Repair-target selection across four pipeline shapes. The causal-root selector is
exact in every cell. The baseline matches it only when declaration order already encodes
execution order.}
\label{tab:selection}
\end{table}

The causal-root selector is exact in all 38 contested cases of every shape and both
orderings. The baseline matches it perfectly when declaration order happens to encode
execution order and falls to 21\,\% when it does not. This gives the graph's value a precise
form: \emph{it is worth exactly the extent to which declaration order diverges from
execution order}. A team whose pipelines are written in order gains nothing here. A team
whose pipelines are model-authored gains the difference between 100 and 21 per cent, and
cannot know in advance which case they are in, which is itself an argument for computing
the relation rather than assuming it.

\subsection{The model judge, and where each signal fails}
\label{sec:compute2}

We expected the language-model judge to be the weak arm. It was not, and reporting what
happened instead is more useful than the result we set out to find.

The judge is a 3B model run locally, asked about one boundary at a time and given that
stage's declared purpose in prose. Across both corpora, 104 silent faults with three
repeats each, it never once disagreed with itself. Every one of the 104 cases came back
unanimous over 312 calls, with no parse failures.

A false-trust rate alone would mean nothing, because a judge answering suspect
indiscriminately scores perfectly while being useless. We therefore asked the same judge
about boundaries known to be sound: every boundary of the clean unfaulted pipeline, and
for each case a boundary strictly upstream of the fault, which cannot be contaminated.
Across 102 such calls it returned trusted every time, so its false-suspicion rate matches
both deterministic signals at zero and the comparison stands.

\begin{table}[t]
\centering
\begin{tabular}{lrrl}
\toprule
Signal & Missed & Rate & 95\,\% CI \\
\midrule
schema validation      & 81/104 & 78\,\% & [69, 85] \\
declared invariants    & 28/104 & 27\,\% & [19, 36] \\
model judge (3B)       & 19/104 & 18\,\% & [12, 27] \\
invariants + judge     &  7/104 &  7\,\% & [3, 13] \\
\bottomrule
\end{tabular}
\caption{Faults left undetected over 104 cases, with Wilson intervals. The intervals for
invariants and the judge overlap; the pair is separated from both.}
\label{tab:signals}
\end{table}

\subsection{What the sample supports, and what it does not}

The ranking in Table~\ref{tab:signals} invites a conclusion the data does not carry, so
we tested it on the paired outcomes. An exact McNemar test does not separate the judge
from declared invariants (12 against 21 discordant, $p=0.16$), and their confidence
intervals overlap. We make no claim that either is the better signal, and the sample is
large enough that this is a finding rather than a shrug: two signals of quite different
kinds, one a compiled check and one a 3B model, are not distinguishable by error rate.

The combination is a different matter. It improves on invariants alone (0 against 21,
$p<0.0001$) and on the judge alone (0 against 12, $p=0.0005$). Both single signals are
significantly worse than using both; neither is significantly worse than the other.
Schema validation is separated from all three at $p<0.0001$, and adding it to any pair
changes nothing, because it catches nothing they do not.

The result is therefore about complementarity rather than ranking. Invariants catch
twelve faults the judge misses, and every one is a matter of counting or shape: fields
removed from a record, sequences silently shortened. The judge catches twenty-one the
invariants miss, and eighteen of those are substituted identifiers, including every
plausible one that defeated the pattern check. Each signal is blind close to where the
other sees, and it is that disjointness the evidence establishes.

It would be wrong to read this as a model outperforming a static check. The judge was
given the stage's declared purpose, for instance that a retrieval stage must fetch every
source in every declared class and report what it retrieved. That is an acceptance
criterion. It is the same kind of statement the invariant encodes, written in prose and
checked by a model rather than compiled into a count and a regular expression.

So the comparison was never between having criteria and not having them. It was between
two encodings of the same criteria, and the encodings fail differently. A count is exact
about quantity and silent about meaning. A sentence carries meaning and is vague about
quantity. That is why the union is so much stronger than either part, and it is a reason
to declare criteria at a boundary at all rather than a reason to prefer one checker.

\subsection{Showing the graph to a model does not pay}
\label{sec:show}

The use an execution graph is most often justified by is presentation: give the agent a
relevant slice rather than everything or nothing. Four regimes were asked the same
question over a branching corpus, name the faulty stage, with the injected fault as the
answer.

\begin{table}[t]
\centering
\begin{tabular}{lrrl}
\toprule
Regime & Accuracy & Context & Note \\
\midrule
stage purposes only        & 15\,\% & 1{,}343 & no captured values \\
size-matched random slice  & 35\,\% & 1{,}402 & control \\
graph-selected region      & 40\,\% & 1{,}466 & $p=0.69$ vs random \\
selected plus expansion    & 46\,\% & 1{,}466 & the deployed method \\
complete current run       & 50\,\% & 1{,}913 & \\
accumulated history, d=64  & 40\,\% & 129{,}246 & \\
\bottomrule
\end{tabular}
\caption{Localisation under six context regimes. Selection is not separable from a random
slice of the same size, including in the form the system actually deploys; accumulated
history costs two orders of magnitude more for less.}
\label{tab:context}
\end{table}

Captured values matter enormously, 15\,\% against 40 to 50 for any regime showing what the
code produced, which is the strongest evidence here that execution records are the right
substrate. Selecting among them does not help: the graph-selected region is not separable
from a size-matched random slice ($p=0.69$), and adding expansion on request, which is the
method as deployed rather than a weakened stand-in, reaches 46\,\% against 50\,\% for
simply showing everything. Accumulated history degrades from 52\,\% at depth four or less
to 40\,\% at depth sixteen or more while its context grows to 129{,}246 characters against
1{,}913. The mechanism is visible in the failures: shown captured values, a model names a
downstream stage in 83\,\% of cases, and selection acts on the prompt rather than on the
answer, so a model handed a selected region still blames the most downstream thing in it.

\subsection{A protocol that withholds and is asked}
\label{sec:collapsed}

The arms above hand the reviewer a slice chosen in advance. The system we study does
something different, and the distinction turns out to matter. Its initial package carries
nodes at exactly the boundary's function-stack level, names the nested helpers without
their contents, and permits the reviewer to request one direct child at a time for at most
three rounds. Evaluating that needs a pipeline whose boundaries call helpers, so we built
one and planted faults at three depths: visible at the boundary, inside a helper, and
inside a sub-helper.

\begin{table}[t]
\centering
\begin{tabular}{lrrrr}
\toprule
Arm & Faults & Clean & Mean chars & Expansions \\
\midrule
collapsed, expansion on request & 5/6 & 4/4 & 841 & 9 \\
everything, flattened           & 5/6 & 4/4 & 922 & 0 \\
random subset, budget matched   & 4/6 & 3/4 & 831 & 0 \\
\bottomrule
\end{tabular}
\caption{The deployed protocol against sending everything and against a size-matched
random subset. It matches the former on nine per cent less context and, unlike every
fixed-slice arm we tested, does not tie the latter.}
\label{tab:collapsed}
\end{table}

This is the only selection arm in the paper that separates from a size-matched random
control, and it forces a narrower statement of the result above. A selected region handed
over in advance is not better than a random one of the same size. A protocol that
withholds nested detail and lets the reviewer pull what it needs is a different mechanism,
and on this evidence it works: the same accuracy as sending everything, for less, and
better than random at equal cost. What we can say is that the benefit lies in the
interaction rather than in the choice of slice, which is not what we expected and not what
the fixed-slice arms could have shown.

Two details qualify it. Faults planted two levels deep were caught with zero and one
expansions, because a deep fault usually surfaces in the boundary's own output, so the
reviewer rarely has to drill to where the fault lives. And the single miss is a reversed
ranking judged clean after an expansion was spent, which is precisely the fault an
ordering scar catches deterministically in Section~\ref{sec:scars}. With ten cases per arm
this is directional rather than sized, and we implemented the collapsed-helper half of the
protocol but not its artifact inspection.

\subsection{External validity: three recovered runs of the real system}
\label{sec:external}

Everything above runs on pipelines we authored, which is the weakest thing about it. Three
records of the real system, deleted in a cleanup commit but recoverable from its git
history, supply the external check.

A four-arm run of 29 July 2026, on the real model-authored pipeline with a frontier model,
reports that graph-selected context matched the full graph's issue coverage with 46\,\%
fewer input tokens on the initial run and 47\,\% on the latest, while accumulated flat
history grew from 71{,}470 to 303{,}757 characters and became 29\,\% more expensive than
the selected region. A second run of 30 July shows the same shape at 31{,}369 tokens for
stage metadata against 109{,}590 for accumulated history, all arms finding the same three
issue functions. This is the efficiency half of our claim, measured on the system itself.

The same 29 July record contains one further observation we regard as the most important
sentence in either dataset: accumulated history \emph{retained a previously repaired
boundary as an open issue after every other arm treated it as resolved}, which its authors
attribute to stale prior versions remaining salient in flat history. Our synthetic history
sweep never produced this, and the recovered run shows why: our accumulated history was
clean prior iterations, which hand the model a helpful baseline to diff against, whereas
real history contains superseded versions of boundaries that have since changed. Staleness,
not volume, is the active ingredient. Section~\ref{sec:stale} tests that as an
intervention.

\subsection{Staleness, briefly}
\label{sec:stale}

Accumulated history harms a model judge, and what it does depends on the model. A
superseded faulty run in context halves one judge's detection of genuine new faults;
another judge cannot recognise a clean run without one, inventing a fault in sixteen of
twenty clean cases and naming the final stage every time; the deployed system's own record
shows a third model resurrecting a resolved issue from flat history. Neither of ours
re-flagged the repaired boundary, which was the predicted failure. The stable conclusion
is a dependence rather than a direction: history content modulates judge behaviour in
model-specific ways that no prompt here controlled, while deterministic signals are
unaffected by context by construction. Full cells in Appendix~\ref{app:stale}; the
detection drop does not survive the multiplicity correction of Section~\ref{sec:power}.

\subsection{Checks placed where faults enter}
\label{sec:scars}

Every signal so far was written in advance by an author deciding what to assert. The
alternative is to learn a check from a fault once it has been caught, place it on the
boundary where it entered, and pay for that fault only once.

We adopt the term \emph{scar} from prior work on coupling kernels, where a scar is a
low-rank invariant subspace that lets a query be \emph{read} exactly rather than searched
over an intractably large state space \cite{rovai2026worldkernel}, and from its use in
agent goal induction, where a goal is treated as a low-rank commitment set read in a
number of probes proportional to its rank \cite{rovai2026civex}. The move transfers
directly. A caught fault commits a small set of properties; rather than re-deriving them
on every subsequent run, they are read off once and pinned to the boundary that produced
them. What is new here is the application to trust decisions and the measurement of what
it buys, which separates the two halves of this paper more sharply than anything else.

A scar is minted from a population of healthy observations of one boundary and one faulty
observation of the same boundary, keeping only structural properties that held across the
whole population and broke in the faulty one: a field present throughout and now absent, a
count below its healthy floor, a list that stopped being homogeneous, an ordering that
held and no longer does. A property expressible as no such check is refused rather than
stored as vague suspicion, because the deterministic signals' zero false-suspicion rate is
the thing that makes them worth having.

\begin{table}[t]
\centering
\begin{tabular}{lrrr}
\toprule
Signal & Detected & Localised & Misplaced \\
\midrule
declared contracts & 26/90 & \textbf{0/90} & 26/90 \\
learned scars      & 60/90 & \textbf{60/90} & 0/90 \\
\bottomrule
\end{tabular}
\caption{Ninety data faults. Localisation is whether the flagged boundary is the one the
fault entered. Declared contracts never localise: every detection fires on an aggregating
stage downstream of the cause.}
\label{tab:scars}
\end{table}

Table~\ref{tab:scars} contains the result we consider most important in this paper.
Hand-written contracts detect 26 of 90 faults and localise \emph{none} of them. A stage
that faithfully passes on a bad input satisfies every criterion anyone thought to write
about its own behaviour, so the checks fire on the aggregating stages downstream and the
suspect set sits entirely past the cause. This is the same failure the model judges
showed, in a signal with no model in it, which tells us the downstream-blame bias is not a
property of language models but of where checks are placed. Scars localise every fault
they detect, because a scar is placed where the fault was observed to enter rather than
where an author's attention happened to fall.

Two further results bound the method. Population size does not affect detection at all,
identical at one, four and eight healthy observations; what it buys is false suspicion,
falling from 21 of 56 boundary-checks at $n=1$ to zero at $n=8$. A single healthy
observation cannot separate an invariant from that run's value, so scars minted from one
pair fire on honest variation, and only a population can tell the difference. Second, a
vocabulary check is minted only when the population shows the domain closed: identifiers
drawn from a fixed registry stop admitting new values, dates and free text do not, and a
vocabulary check over an open domain is guaranteed to fire on an honest run eventually.
Enforcing that costs detection, 73 falling to 60, and takes false suspicion on held-out
healthy runs the scars never saw from 3 of 28 to zero. We regard that trade as the correct
one, and the refusal of faults visible only in an open domain as a feature of the design.

We also implemented the two graph computations this suggests, and both are negative.
Forward propagation of contamination reclassifies nothing, because the declared checks
fire downstream of the fault and there is nothing to propagate forward from. Backward
attribution narrows to seven boundaries of seven, because in a pipeline that funnels
through a merge every branch is an ancestor of the first suspect. Both failures have one
cause: exoneration requires a check that the fault \emph{would have violated}, not merely
a check that passed, and a hand-written contract does not supply one where a scar does.

\section{Defects, in the system and in the measurement}
\label{sec:defects}

\subsection{In the trust machinery}

Four defects, detailed in Appendix~\ref{app:defects}: the repair-scope validator silently
skipped enforcement whenever a declared boundary name failed to resolve, a branch the
authoring model controls and that the entire workflow test suite unknowingly ran through;
the module producing the published numbers had no tests; resource limits around generated
code were inactive by default; and the real home directory, with credentials, reached
model-written code that has network access outside replay
\cite{greshake2023injection}.

\subsection{In our own measurement}

A benchmark built to verify agents has no standing to measure anyone else if it cannot
catch its own defects, and five times a confident number here turned out to be an artifact
of how we built the experiment.

A control arm shown the source text of the implicated region scored 75\,\%, appearing to
refute the premise that captured values matter; our faults were source mutations, so it
was reading the fault rather than diagnosing it, and moving faults into recorded responses
with the source byte-identical across all 48 cases dropped it to 0\,\%. A 3B judge scored
0\,\% on seven-way localisation while returning valid stage names, answering one of two
regardless of input; a 30B model later returned \texttt{transform\_3} for all 48 cases at
every chain length from five to seventeen stages, producing a flat apparent degradation
with horizon that was a constant. Both look like results until the answer distribution is
examined, and neither can exhibit an effect of context because neither depends on context.
We reported that graph-selected context beat a size-matched random slice, 40 against
28\,\%, on a third of the eventual sample; at full sample the difference vanished
($p=0.69$), and the claim was retracted, having been made in the same paragraph as a
criticism of single-run reporting. Because data-level faults land on the stage that
fetches the corrupted response, ground truth concentrated on three of seven boundaries,
making the correct chance baseline 42\,\% rather than 14\,\%, and an arm shown no captured
values scored exactly that by always naming the most common answer. Finally, a dump cap truncated retrieval
evidence before the judge saw it, so faults on later sources were invisible and the judge
blamed downstream aggregates. Each would have survived into a table, and we take the
frequency as evidence that scoring against constructed ground truth is worth as much for
what it invalidates as for what it establishes.

\section{Limitations}

Three judges appear here, a 3B, a 30B mixture and one frontier model, with large and
non-uniform capability differences, so no rate should be read as a property of models in
general; and the judge was never asked to choose among boundaries the way the selector
does, so the selection and detection results are not one story.

Faults are injected rather than arising from a model, and model-authored faults would
differ systematically from ours; scars inherit this doubly, learning only differences
observable at the boundary where the fault entered, and evaluated on one fault taxonomy.
Our contamination model follows module-level dataflow, so state passed through a file, a
global or a side effect defeats it, which makes it sound for conforming pipelines and
unsound for the ones most likely broken. The recovered records corroborate the efficiency claim without
bearing on detection rates, and no pipeline exceeds seventeen stages, so nothing here
measures the true long-horizon regime \cite{chen2026horizon}.


\bibliographystyle{plain}
\bibliography{refs}

\appendix

\section*{Artifact availability}
The benchmark, every corpus, all result files and the analysis scripts accompany this
submission as an anonymized artifact and will be released publicly on acceptance; the
recovered records of Section~\ref{sec:external} are identified by the commit hash of the
repository they were recovered from.

\section{The published comparison}
\label{app:published}

\begin{table}[h]
\centering
\begin{tabular}{lrrrr}
\toprule
Arm & Repairs & Effective & Input tokens & Unresolved boundaries \\
\midrule
\texttt{semantic\_only}  & 3 & 0 & 568{,}715   & 4 \\
\texttt{history\_full}   & 3 & 1 & 1{,}824{,}533 & 4 \\
\texttt{etiq\_full}      & 3 & 0 & 1{,}331{,}686 & 3 \\
\texttt{etiq\_selected}  & 3 & 1 & 1{,}000{,}797 & 2 \\
\bottomrule
\end{tabular}
\caption{The published comparison, one run per arm, 30 July 2026, taken from the
archived machine-readable summary. Its own record notes that ``no arm reached final
trust''.}
\label{tab:published}
\end{table}

\section{Defects in the trust machinery, in detail}
\label{app:defects}

The validator confining a repair to one boundary returned early whenever the target could
not be resolved to a function, so in that branch a repair could rewrite the whole target
file. The branch is reached whenever a declared boundary name does not resolve in the
generated source, and the authoring model writes both, so it is reachable by ordinary
error and steerable on purpose; the harness splices over the same span, so an unresolvable
target produced a whole-file substitution the validator accepted. It was untested, because
the scanner fixture reported a call stack inside functions the authoring fixture never
defined, so no target ever resolved and every workflow test of repair ran through the
unenforced branch. Separately, the module producing the published numbers was imported by
no test, resource limits around generated code defaulted to inactive, and the home
directory reached the generated pipeline, which has network access outside replay mode;
captured third-party text now reaches the reviewer inside a digest-keyed fence it cannot
close \cite{greshake2023injection}.

\section{A declared invariant, concretely}
\label{app:invariant}

The retrieval stage of the branching pipeline carries: \texttt{retrieve\_registries} must
return \texttt{source\_ids} and \texttt{source\_class}; \texttt{source\_ids} must hold at
least three entries; every entry must match \texttt{\^{}[a-z]+\_[a-z]+\$}. Three clauses
catching a disappeared field, a silently shortened retrieval, and a malformed identifier;
what none can catch is an identifier that is correctly shaped and simply wrong, which is
measured rather than assumed in the body.

\section{Extended related work}
\label{app:related}

Empirical work on why multi-agent systems fail supplies the closest framing for our
problem: Cemri et al.\ build a taxonomy from more than two hundred annotated traces and
make task verification one of its three top-level categories \cite{cemri2025multiagent},
though their failures are annotated by human experts rather than detected by the system at
runtime; Chen et al.\ survey the same territory as a gap between what long-horizon agents
are asked to do and what their planning, memory and evaluation machinery supports
\cite{chen2026horizon}. The assumption our benchmark tests is that a language model
supplies that verification, and there is reason for caution: judges favour text more
familiar to themselves \cite{wataoka2024selfpreference}, which matters because reviewer
and reviewed here share a model family; a judge's self-agreement can be high while its
agreement with the truth is not \cite{norman2026reliability}, the distinction our
localisation metric operationalises; verifiers checking only what they can express admit
confident false positives \cite{helff2026gaming}; and multi-step agents exploit
naturalistic chances to tamper with what evaluates them \cite{thaman2026rewardhacking}.
Provenance-based trust for agents \cite{wang2026traces} and recoverable execution from
recorded checkpoints \cite{zhuang2026agentrewind} are the constructive lines this work
measures rather than proposes.

Our method borrows its central move from mutation testing, a defect of known location
measuring whether a suite detects it, including the equivalent-mutant problem we meet as
an injection that changes nothing \cite{jia2011mutation}. The substrate is execution
provenance \cite{freire2008provenance}, whose distinction between a record of what a
computation did and a claim about what it meant locates this paper's difficulty: the graph
is a record, the trust label is a claim, and only the first is checkable. The fencing of
Section~\ref{sec:defects} addresses indirect prompt injection through fetched content
\cite{greshake2023injection}, a designed-in exposure here since the pipeline must retrieve
third-party sources at runtime.

\section{Staleness intervention, full results}
\label{app:stale}


We crossed two factors, twenty cases per cell on each of two judges: the history shown
(clean prior iterations, or a superseded run carrying a fault since repaired) and the
current run (clean, or carrying a new fault at a different boundary), asking for the
faulty stage of the current run or none.

\begin{table}[h]
\centering
\begin{tabular}{llrr}
\toprule
History & Current run & Judge A & Judge B \\
\midrule
clean & clean     & 18/20 & 4/20 \\
stale & clean     & 17/20 & 13/20 \\
clean & new fault & 18/20 & 19/20 \\
stale & new fault & 11/20 & 19/20 \\
\bottomrule
\end{tabular}
\caption{Correct decisions under the staleness intervention. Each judge fails in a
different place: A dismisses genuine new faults when a stale prior is present, B cannot
recognise a clean run unless one is.}
\label{tab:stale}
\end{table}

Neither judge re-flagged the repaired boundary, which was the predicted failure; that
happened once in 160 stale-history calls. Judge A, handed a superseded faulty prior,
misses genuine new faults, 18/20 falling to 11/20 (Fisher exact, $p=0.03$ uncorrected,
suggestive after the multiplicity correction of Section~\ref{sec:power}); the prior appears
to act as a contrast anchor, and a current run that looks better than the old one reads as
fine. Judge B shows the opposite defect: it detects the new fault regardless of history,
19/20 in both cells, but cannot recognise a clean run, inventing a fault in 16 of 20
clean-history cases and naming the final stage in every one, which is the downstream-blame
bias with nothing to blame; a superseded faulty prior partially restores its ability to
answer clean, 13/20 against 4/20. The deployed system's own record supplies a third
signature, a frontier model resurrecting a resolved issue from flat history.

Three judges, three failure shapes from one manipulation. The stable conclusion is a
dependence rather than a direction: the content of accumulated history modulates a model
judge's trust decisions in model-specific ways that no prompt here controlled, while
deterministic signals are unaffected by context by construction.


\end{document}
